% !TeX program = lualatex
% !TeX encoding = utf8
% !TeX spellcheck = uk_UA
% !TeX root = ../main.tex
% ElMagConspect.tex



%=========================================================
\Opensolutionfile{answer}[\currfilebase/\currfilebase-Answers]
\Writetofile{answer}{\protect\section*{\nameref*{\currfilebase}}}
\chapter{Енергія електричного поля}\label{\currfilebase}\hypertarget{ElEnergy}{}
\graphicspath{{\currfilebase/Pictures}}
%=========================================================
\epigraph{\Alexander Говорячи про енергію поля, я хочу, щоб мене розуміли буквально.}{Джеймс
Клерк Максвелл, \enquote{Динамічна теорія електромагнітного поля} (1865)}

У~попередніх розділах ми розібрали, як поводиться електричне поле у~вакуумі, у~провідниках
і~в~діелектриках. Ми знаємо, як обчислити поле за заданим розподілом зарядів (теорема
Гаусса, принцип суперпозиції), і~як поводиться поле на межі розділу середовищ. Тепер з'ясуємо,
\emph{скільки енергії зберігає електричне поле} і~\emph{як за цією енергією обчислювати
сили}, що діють на заряджені тіла.

Цей розділ складається з~трьох частин. Спершу ми введемо \emphx{електричну ємність}\index{Ємність електрична}~---
величину, яка описує, скільки заряду (а~отже, і~поля) здатна накопичити система
провідників при заданій різниці потенціалів. Далі ми покажемо, що енергія заряджених
тіл~--- це насправді енергія \emph{поля}, розподілена в~просторі з~густиною $w =
\Efield\cdot\Dfield/8\pi$. Нарешті, ми навчимось обчислювати \emphx{пондеромоторні сили}
--- сили, що діють на заряджені провідники й~діелектрики в~електричному полі,~--- за
допомогою методу віртуальних переміщень.


%% --------------------------------------------------------
\section{Електрична ємність}\hypertarget{electric-capacitance}{}
%% --------------------------------------------------------


%%% --------------------------------------------------------
\subsection*{Ємність відокремленого провідника}
%%% --------------------------------------------------------


Нехай відокремлений (далекий від інших тіл) провідник несе заряд $q$; тоді, як ми знаємо,
весь цей заряд розподіляється по його поверхні, а~сам провідник~--- еквіпотенціальне тіло
з~деяким потенціалом $\phi$. Якщо заряд збільшити в~$n$ разів, то, за принципом
суперпозиції, поле в~кожній точці простору також зросте в~$n$ разів, а~отже, в~$n$ разів
зросте й~робота з~переміщення пробного заряду з~поверхні провідника на нескінченність,
тобто й сам потенціал $\phi$. Це означає, що відношення
\begin{equation}\label{eq:CapacitorDef}
	\tcbhighmath{
		C = \frac{q}{\phi}
	}
\end{equation}
не залежить від заряду провідника і~визначається лише його формою, розмірами та
властивостями середовища, що його оточує. Цю величину називають \emphx{електричною ємністю
провідника}.

Зокрема, ємність не залежить від матеріалу провідника. Означення
\eqref{eq:CapacitorDef} має сенс лише для \emph{відокремленого} провідника, коли поблизу
немає інших тіл, здатних помітно змінити його потенціал.



%% --------------------------------------------------------
\subsection*{Два провідники: взаємна ємність}
%% --------------------------------------------------------


Найпоширеніший спосіб зарядити провідник~--- перенести частину заряду з~другого: тоді на
одному з'являється заряд \(+q\), на іншому \(-q\). Для такого стану вводять характеристику пари~---
\begin{equation}\label{eq:MutualCapacitanceDef}
	\tcbhighmath{
		C = \frac{q}{\phi_+ - \phi_-}
	}
\end{equation}
де \(\phi_\pm\)~--- потенціали провідників із зарядами \(\pm q\).
Величину~\eqref{eq:MutualCapacitanceDef} називають \emphx{взаємною ємністю}\index{Ємність електрична!взаємна} пари.

%% --------------------------------------------------------
\subsection*{Коефіцієнти ємності*}
%% --------------------------------------------------------
\begin{Extra}
Для\index{Ємність електрична!коефіцієнти ємності} відокремленого провідника заряд і~потенціал пропорційні: \(q = C\phi\). Якщо провідників
кілька, потенціал кожного залежить від зарядів усіх сусідів (принцип суперпозиції), тому
\begin{equation}\label{eq:Cij}
	q_i = \sum_j C_{ij}\,\phi_j.
\end{equation}
Коефіцієнти \(C_{ij}\) залежать лише від форми, розмірів і~взаємного розташування
провідників (і~від проникності середовища), але не від самих зарядів чи потенціалів.

Діагональні \(C_{ii}\)~--- \emphx{власні ємності}\index{Ємність електрична!власна}: заряд провідника \(i\) при одиничному
власному потенціалі й~заземлених сусідах. Недіагональні \(C_{ij}\) (\(i\neq j\))~---
\emph{коефіцієнти взаємного впливу}: заряд, наведений на заземленому провіднику \(i\),
коли сусід \(j\) піднято до одиничного потенціалу.

Знаки коефіцієнтів визначені однозначно: \(C_{ii}>0\), а~\(C_{ij}\le 0\) при \(i\neq j\),
бо, піднімаючи потенціал одного провідника, ми завжди наводимо на заземлених сусідах заряд
\emph{протилежного} знака. Крім того, \(C_{ij}=C_{ji}\)~--- наслідок теореми взаємності
електростатики (її доведення виходить за межі курсу). Фізично це означає: байдуже, який
із двох провідників підняти до потенціалу, а~який заземлити,~--- наведений на сусідові заряд
буде той самий. У~прикладі з~двома кулями нижче \(C_{12}\approx -C_1C_2/b\) справді
однакове за будь-якого вибору.

\emph{Оцінка зверху для \(C_{12}\).} Розглянемо пару провідників з~однаковою власною
ємністю \(C_{11}=C_{22}\equiv C_0\) (наприклад, дві однакові кулі). Піднімемо потенціал
провідника~1 до одиниці, а~провідник~2 заземлимо (\(\phi_2=0\)). Тоді, згідно
з~\eqref{eq:Cij}, заряд першого провідника дорівнює \(q_1=C_{11}=C_0\), а~на другому
наводиться заряд \(q_2=C_{12}\) (від'ємний, за щойно сформульованим правилом знаків).

	\begin{SCfigure}[0.75][h!]\centering
		\localinput{field-lines-bound.tikz}
		\caption{Частина силових ліній провідника~1 замикається на заземленому
		провіднику~2, решта йде на нескінченність\label{pic:FieldLinesBound}}
	\end{SCfigure}

Фізично \(q_2\)~--- це заряд, \enquote{перехоплений} тими силовими лініями поля
провідника~1, що замикаються саме на заземленому сусідові
(рис.~\ref{pic:FieldLinesBound}); решта ліній іде на нескінченність (або на інші тіла).
Оскільки сусід не може перехопити більше силових ліній, ніж їх випустив сам провідник~1,
наведений заряд не може за модулем перевищити початковий:
\begin{equation}\label{eq:OffDiagonalBound}
	\tcbhighmath{
		|C_{12}| \le C_0.
	}
\end{equation}
Рівність \(C_{12}\to -C_0\) означає, що \emph{всі} силові лінії від провідника~1
замикаються на провіднику~2~--- жодна не йде на нескінченність, тобто екранування повне.
Саме до цієї межі прямує тісно зближена пара провідників~--- \emphx{конденсатор}, з~яким
ми детальніше познайомимось нижче.

Щоб виразити \(C\) через коефіцієнти \(C_{ij}\), підставимо \(q_1=+q\), \(q_2=-q\)
у~систему~\eqref{eq:Cij} (з~урахуванням симетрії \(C_{21}=C_{12}\)):
\begin{equation*}
	q = C_{11}\phi_1 + C_{12}\phi_2, \qquad
	-q = C_{12}\phi_1 + C_{22}\phi_2.
\end{equation*}
Розв'язавши цю лінійну систему відносно \(\phi_1\), \(\phi_2\) (звичайним оберненням
\(2\times2\)-матриці), знаходимо, що різниця потенціалів пропорційна заряду:
\begin{equation*}
	\phi_1 - \phi_2 = \frac{q\,(C_{11}+C_{22}+2C_{12})}{C_{11}C_{22}-C_{12}^2},
\end{equation*}
звідки за означенням~\eqref{eq:MutualCapacitanceDef}
\begin{equation}\label{eq:MutualCapacitanceGeneral}
	\tcbhighmath{
		C = \frac{C_{11}C_{22}-C_{12}^2}{C_{11}+C_{22}+2C_{12}}.
	}
\end{equation}

Для \emph{симетричної пари} \(C_{11}=C_{22}\equiv C_0\) чисельник розкладається як
різниця квадратів, \(C_0^2-C_{12}^2=(C_0-C_{12})(C_0+C_{12})\), а~знаменник дорівнює
\(2(C_0+C_{12})\); після скорочення спільного множника \((C_0+C_{12})\):
\begin{equation}\label{eq:SymmetricPairCap}
	\tcbhighmath{
		C = \frac{C_0 - C_{12}}{2}.
	}
\end{equation}

Ця формула добре узгоджується з~оцінкою~\eqref{eq:OffDiagonalBound} і~показує, як взаємна
ємність зростає в~міру екранування: якщо взаємний вплив зникає (\(C_{12}\to 0\),
провідники далеко один від одного), то \(C\to C_0/2\)~--- це просто послідовне з'єднання
двох однакових ємностей \(C_0\); якщо ж екранування майже повне (\(C_{12}\to -C_0\),
гранично допустиме за~\eqref{eq:OffDiagonalBound}), то \(C\to C_0\)~--- поле майже все
замкнене між обкладками. Саме другий режим реалізується в~конденсаторі.

\emph{Приклад: дві далекі кулі.} Дві кулі на відстані \(b\), значно більшій за їхні
розміри (рис.~\ref{pic:TwoSpheres}). Нехай \(C_1\), \(C_2\)~--- ємності куль поодинці
(\(C_{11}\to C_1\), \(C_{22}\to C_2\) при \(b\to\infty\)).

	\begin{SCfigure}[1.5][h!]\centering
		\localinput{two-spheres.tikz}
		\caption{Дві далекі одна від одної заряджені кулі\label{pic:TwoSpheres}}
	\end{SCfigure}

	Зарядимо кулі \(\pm q\) і~врахуємо вплив сусідньої кулі як точкового заряду на
	відстані \(b\):
	\begin{equation*}
		\phi_+ = \frac{q}{C_1} - \frac{q}{b}, \qquad
		\phi_- = -\frac{q}{C_2} + \frac{q}{b}.
	\end{equation*}
	Звідси \(\phi_+-\phi_- = q\left(\dfrac1{C_1}+\dfrac1{C_2}-\dfrac2b\right)\)
	(поляризацією куль, що дає поправки порядку \(1/b^2\), знехтувано), і~з~\(1/(1-x)\approx
	1+x\) при малому \(x = 2C_1C_2/(C_1+C_2)b\)
	\begin{equation}\label{eq:TwoSpheresCap}
		\tcbhighmath{
			C \approx \frac{C_1C_2}{C_1+C_2}\left(1 + \frac{C_1C_2}{C_1+C_2}\frac{2}{b}\right).
		}
	\end{equation}
	У~головному наближенні (\(b\to\infty\)) взаємна ємність дорівнює ємності послідовного
	з'єднання; поправка описує слабкий взаємний вплив куль.

	Той самий результат дає загальна формула~\eqref{eq:MutualCapacitanceGeneral}. Якщо
	кулю~2 заземлити, на ній наводиться заряд \(q_2\approx -C_2q_1/b\), що компенсує
	потенціал заряду \(q_1\), тобто \(C_{12}\approx -C_1C_2/b<0\). Підставивши
	\(C_{11}\approx C_1\), \(C_{22}\approx C_2\) і~цей \(C_{12}\), з~точністю до \(1/b\)
	знову отримуємо~\eqref{eq:TwoSpheresCap}; з~неправильним знаком \(C_{12}\) поправка
	вийшла б протилежною.

Наскільки повним може бути це екранування, видно з~\eqref{eq:SymmetricPairCap}: коли
обкладки зближуються, коефіцієнт взаємного впливу прямує до своєї межі \(C_{12}\to -C_0\)
(оцінка~\eqref{eq:OffDiagonalBound}), а~отже,
\begin{equation*}
	C \to C_0,
\end{equation*}
тобто взаємна ємність прямує до власної. Не слід розуміти це так, ніби ємність
конденсатора дорівнює ємності однієї відокремленої пластини: \(C_0 = C_{11}\)~--- це
ємність пластини \emph{поряд із заземленою сусідньою}, і~при зближенні пластин вона сама
зростає (для плоского конденсатора~--- до \(\epsilon S/4\pi d\), див. нижче). У~цій
границі поле майже все замкнене між обкладками, і~\(C\) перестає залежати від того, що
відбувається далеко.
\end{Extra}



%% --------------------------------------------------------
\subsection*{Конденсатор}
%% --------------------------------------------------------



\emphx{Конденсатор}\index{Конденсатор}~--- система з~двох провідників (\emph{обкладок}), розташованих так
близько (зазвичай через тонкий шар діелектрика), що майже все поле зосереджене у~проміжку
між ними. Завдяки такому \emph{екрануванню} взаємна ємність конденсатора, на відміну від
ємності відокремленого провідника чи пари віддалених провідників, практично не залежить від
інших тіл поблизу~--- саме тому вона є зручною й~відтворюваною характеристикою.
Формула~\eqref{eq:MutualCapacitanceDef} лишається тією самою; конденсатор~--- окремий
випадок пари провідників з~екранованим полем.

Чому поле конденсатора майже повністю зосереджене між обкладками і~чому його ємність
більша, ніж в~окремих обкладок, пояснено вище, у~підрозділі \enquote{Коефіцієнти
ємності*}.

\begin{Example}{Обкладки з~різними зарядами}
	Нехай на обкладках плоского вакуумного конденсатора (площею \(S\) кожна)~--- не
	протилежні, а~\emph{довільні} заряди \(q_1\) і~\(q_2\) (наприклад, пластини зарядили
	окремо, а~потім розташували паралельно на малій відстані одна від одної).

    Розкладемо їх на \emph{антисиметричну} і~\emph{симетричну} частини:
	\begin{equation*}
		q_a = \frac{q_1-q_2}{2}, \qquad q_s = \frac{q_1+q_2}{2}, \qquad q_1 = q_a+q_s,\quad q_2=-q_a+q_s.
	\end{equation*}
	Заряд \(q_a\) поводиться як звичайний заряд конденсатора~--- сідає на
	\emph{внутрішні} (обернені одна до одної) поверхні обкладок і~створює однорідне поле
	лише в~проміжку між ними, \(E_\text{між}=4\pi q_a/S\). Заряд \(q_s\), навпаки,
	розподіляється порівну на \emph{зовнішніх} поверхнях: поле від двох однаково
	заряджених паралельних площин між ними взаємно компенсується (теорема Гаусса), тому
	в~проміжку \(q_s\) поля не дає, а~весь цей заряд виявляється \emph{назовні}, створюючи
	там поле \(E_\text{зовні}=4\pi q_s/S\)~--- так само, як біля поверхні звичайного
	зарядженого провідника. Підставляючи означення \(q_a,q_s\):
	\begin{equation}\label{eq:UnequalPlatesField}
			E_\text{між} = \frac{2\pi(q_1-q_2)}{S}, \qquad E_\text{зовні} = \frac{2\pi(q_1+q_2)}{S}.
	\end{equation}
	Якщо \(q_1=-q_2=q\) (звичайний конденсатор), то \(E_\text{зовні}=0\)~--- поле повністю
	екрановане, як і~слід з~означення конденсатора. Якщо ж, навпаки, \(q_1=q_2\) (обкладки
	заряджені однаково), то \(E_\text{між}=0\), а~весь заряд \enquote{виходить}
	назовні~--- пара обкладок при цьому поводиться просто як два майже злиті окремі
	провідники.
\end{Example}

Повне екранування~--- ідеалізація. Частина силових ліній реальних обкладок завжди йде на
треті тіла (землю, корпус приладу, сусідні провідники), і~тоді заряд обкладки залежить не
лише від різниці потенціалів між обкладками, а~й~від потенціалів сусідів. Таку систему
називають \emphx{складним конденсатором}\index{Конденсатор!складний}; її описує загальна
формула~\eqref{eq:Cij}.


%% --------------------------------------------------------
\subsection*{З'єднання конденсаторів}
%% --------------------------------------------------------


На\index{Конденсатор!з'єднання} практиці конденсатори часто з'єднують у~батареї~--- послідовно або паралельно,~--- щоб
отримати потрібну сумарну ємність.

\begin{figure}[h!]
	\centering
	% --- (а) послідовне з'єднання ---
	\begin{subfigure}[t]{0.48\linewidth}
		\centering
		\resizebox{\linewidth}{!}{\localinput{CapacitorsSeries.tikz}}
		\caption{Послідовне з'єднання}
		\label{pic:CapacitorsSeries}
	\end{subfigure}\hfill
	% --- (б) паралельне з'єднання ---
	\begin{subfigure}[t]{0.48\linewidth}
		\centering
		\resizebox{\linewidth}{!}{\localinput{CapacitorsParallel.tikz}}
		\caption{Паралельне з'єднання}
		\label{pic:CapacitorsParallel}
	\end{subfigure}
	\caption{З'єднання конденсаторів: (а) послідовне~--- заряд на всіх обкладках
	однаковий, напруги додаються; (б) паралельне~--- напруга на всіх конденсаторах
	однакова, заряди додаються}
	\label{pic:CapacitorsConnection}
\end{figure}

\emphx{Послідовне з'єднання.} Нехай $n$ конденсаторів з'єднано послідовно
(рис.~\ref{pic:CapacitorsSeries}). Оскільки обкладки сусідніх конденсаторів, з'єднані між
собою провідником, електрично ізольовані від решти кола, повний заряд кожної внутрішньої
пари обкладок мусить залишатися таким, яким був спочатку, тобто нульовим (конденсатори
до з'єднання не заряджені): якщо на крайню ліву обкладку батареї нанести
заряд $+q$, то шляхом електростатичної індукції той самий за модулем заряд $q$ з'явиться
на \emph{кожному} конденсаторі батареї. Загальна різниця потенціалів (напруга) на батареї
дорівнює сумі напруг на окремих конденсаторах:
\begin{equation*}
	U = U_1 + U_2 + \ldots + U_n = \frac{q}{C_1} + \frac{q}{C_2} + \ldots + \frac{q}{C_n}.
\end{equation*}
Оскільки ємність батареї за означенням $C = q/U$, маємо
\begin{equation}\label{eq:SeriesCap}
	\tcbhighmath{
		\frac1C = \sum\limits_{i=1}^n \frac1{C_i}.
	}
\end{equation}

\emphx{Паралельне з'єднання.} Нехай тепер $n$ конденсаторів з'єднано паралельно
(рис.~\ref{pic:CapacitorsParallel}): усі ліві обкладки з'єднані між собою і~мають спільний
потенціал, так само як і~всі праві. Тому напруга на кожному конденсаторі однакова
і~дорівнює загальній напрузі батареї $U$, а заряди на окремих конденсаторах додаються:
\begin{equation*}
	q = q_1 + q_2 + \ldots + q_n = C_1 U + C_2 U + \ldots + C_n U.
\end{equation*}
Звідси
\begin{equation}\label{eq:ParallelCap}
	\tcbhighmath{
		C = \sum\limits_{i=1}^n C_i.
	}
\end{equation}

Отже, послідовне з'єднання завжди \emph{зменшує} ємність (вона менша за найменшу з~$C_i$),
а~паралельне~--- \emph{збільшує}. Формули такі самі, як для опорів, але \enquote{навпаки},
бо роль опору тут відіграє $1/C$: напруга на конденсаторі $U = q/C$.


%%% --------------------------------------------------------
\subsection*{Обчислення ємності найпростіших конденсаторів}
%%% --------------------------------------------------------




Усі задачі на обчислення ємності розв'язуються за \emph{єдиним алгоритмом}:

\begin{enumerate}
	\item Задати заряд $q$ на обкладках і~знайти розподіл поля $\Efield$ між ними (за
	  теоремою Гаусса, з~урахуванням діелектрика, якщо він є).
	\item Обчислити різницю потенціалів $\Delta\phi = \phi_+ - \phi_- = \int \Efield\cdot
	  d\vect{r}$ уздовж силової лінії.
	\item За означенням~\eqref{eq:MutualCapacitanceDef} знайти $C = q/\Delta\phi$.
\end{enumerate}

Продемонструємо цей алгоритм на кількох типових прикладах. Скрізь вважатимемо, що простір
між обкладками заповнений однорідним діелектриком із проникністю $\epsilon$ (результати
про поле в~діелектрику~--- з~розділу~\ref{FieldInDielectrics}).

\emph{Ємність відокремленої кулі.} Провідна куля радіусом $R$, що несе заряд $Q$, створює
поза собою поле точкового заряду, а~її власний потенціал дорівнює потенціалу на поверхні,
$\phi = Q/(\epsilon R)$. Звідси відразу
\begin{equation}\label{eq:SphereCap}
	\tcbhighmath{
		C = \epsilon R.
	}
\end{equation}
Цей простий результат корисно пам'ятати як орієнтир: ємність провідного тіла має порядок
величини його лінійних розмірів (помножених на $\epsilon$).

\emph{Плоский конденсатор.} Нехай обкладки~--- дві паралельні пластини площею $S$ кожна,
розташовані на відстані $d \ll \sqrt{S}$ одна від одної (щоб знехтувати крайовими
ефектами), а~простір між ними заповнений діелектриком з~проникністю $\epsilon$. Поле між
обкладками однорідне, $E = 4\pi\sigma/\epsilon = 4\pi q /(\epsilon S)$, тому різниця
потенціалів $\Delta\phi = Ed = 4\pi qd/(\epsilon S)$, і
\begin{equation}\label{eq:FlatCap}
	\tcbhighmath{
		C = \frac{\epsilon S}{4\pi d}.
	}
\end{equation}

\emph{Сферичний конденсатор.} Обкладки~--- дві концентричні сфери радіусів $R_1 < R_2$,
простір між якими заповнений діелектриком з~проникністю $\epsilon$. За теоремою Гаусса
поле між обкладками таке саме, як і~поле точкового заряду $q$ у~діелектрику, $E =
q/(\epsilon r^2)$. Різниця потенціалів
\begin{equation*}
	\Delta\phi = \int\limits_{R_1}^{R_2} E\, dr = \frac{q}{\epsilon}\left(\frac1{R_1} - \frac1{R_2}\right) = \frac{q}{\epsilon}\frac{R_2 - R_1}{R_1
		R_2},
\end{equation*}
звідки
\begin{equation}\label{eq:SphericalCap}
	\tcbhighmath{
		C = \frac{\epsilon R_1 R_2}{R_2 - R_1}.
	}
\end{equation}
Легко переконатись, що при $R_2 \to \infty$ (друга обкладка йде на нескінченність)
формула~\eqref{eq:SphericalCap} переходить у формулу~\eqref{eq:SphereCap} для відокремленої
кулі, як і~повинно бути.

\emph{Циліндричний конденсатор.} Обкладки~--- дві коаксіальні циліндричні поверхні
радіусів $R_1 < R_2$ і~довжини $\ell \gg R_2$ (щоб знехтувати крайовими ефектами), простір
між ними заповнений діелектриком з~проникністю $\epsilon$. За теоремою Гаусса поле між
обкладками $E = \dfrac{2q}{\epsilon \ell r}$ (поле нескінченної зарядженої нитки
в~діелектрику), тому
\begin{equation*}
	\Delta\phi = \int\limits_{R_1}^{R_2} E\, dr = \frac{2q}{\epsilon\ell}\ln\frac{R_2}{R_1},
\end{equation*}
звідки
\begin{equation}\label{eq:CylindricalCap}
	\tcbhighmath{
		C = \frac{\epsilon \ell}{2\ln\left(R_2/R_1\right)}.
	}
\end{equation}

Усі чотири формули мають вигляд $C=\epsilon\,\Lambda$, де $\Lambda$~--- величина
розмірності \emph{довжини}, що залежить лише від геометрії ($R$; $S/4\pi d$;
$R_1R_2/(R_2-R_1)$; $\ell/2\ln(R_2/R_1)$). Тому в~гауссовій системі ємність вимірюють
у~сантиметрах ($1$~см $\approx 1{,}11$~пФ). Запам'ятовувати формули окремо не потрібно:
кожну легко вивести за алгоритмом вище.

%\begin{Example}{Числовий приклад}
%	Розглянемо плоский конденсатор із площею обкладок $S = 100\ \text{см}^2 = 10^{-2}\ \text{м}^2$ і відстанню між ними
%	$d = 1\ \text{мм} = 10^{-3}\ \text{м}$, заповнений повітрям ($\epsilon \approx 1$). У гауссовій системі одиниць ємність такого конденсатора
%	$C = \epsilon S/(4\pi d) = 10^{-2}/(4\pi\cdot 10^{-3})\ \text{см} \approx 0{,}8\ \text{см}$. У практичних одиницях (СІ) $1\ \text{см}$ ємності
%	відповідає приблизно $1{,}1\ \text{пФ}$, тож $C \approx 0{,}9\ \text{пФ}$. Якщо між обкладками вставити слюду з $\epsilon \approx 6$, ємність
%	зросте до $\approx 5\ \text{пФ}$ --- саме так \enquote{працює} діелектрик у конденсаторі: не змінюючи геометрії, він у $\epsilon$ разів збільшує
%	ємність.
%\end{Example}

Для порівняння: ємність Землі як відокремленої кулі $C = R_\oplus \approx 6{,}4\cdot10^8$~см
$\approx 710$~мкФ, а~\emph{суперконденсатори} мають ємність у~тисячі фарад~--- у~мільйони
разів більшу. Її досягають завдяки надрозвиненій поверхні пористих вуглецевих електродів
($1000$--$2600$~м$^2$/г) і~\emph{подвійному електричному шару}, в~якому іони електроліту
розділені відстанню, меншою за $1$~нм; обидва чинники входять у~формулу плоского
конденсатора $C = \epsilon S/4\pi d$. Проте за питомою енергією суперконденсатори
($\sim5$--$10$~Вт$\cdot$год/кг) значно поступаються літій-іонним акумуляторам
($\sim100$--$250$~Вт$\cdot$год/кг, енергія в~хімічних зв'язках, а~не в~полі), не кажучи вже
про бензин ($\approx12\,000$~Вт$\cdot$год/кг). Тому вони доповнюють акумулятори там, де
важлива питома \emph{потужність}: рекуперативне гальмування, стабілізація напруги, пікові
навантаження.



%% --------------------------------------------------------
\section{Енергія електричного поля}\hypertarget{field-energy}{}
%% --------------------------------------------------------


%%% --------------------------------------------------------
\subsection*{Взаємна енергія системи точкових зарядів}
%%% --------------------------------------------------------


Щоб зблизити два точкових заряди $q_1$ і~$q_2$ з~нескінченності до відстані $r_{12}$ один
від одного, потрібно виконати проти сил поля роботу $A = q_1q_2/r_{12}$. Цю роботу
і~ототожнюють з~\emphx{взаємною енергією}\index{Енергія зарядів!взаємна} пари зарядів: $U_{12} = q_1q_2/r_{12}$. Для
системи з~$n$ точкових зарядів (рис.~\ref{pic:ChargesTriangle}) сумарна взаємна енергія
набирається з~енергій усіх можливих пар:
\begin{equation}\label{eq:PairEnergy}
	\tcbhighmath{
		W = \frac12 \sum\limits_{i=1}^n \sum\limits_{\substack{j=1\\ j\neq i}}^n \frac{q_iq_j}{r_{ij}},
	}
\end{equation}
де множник $1/2$ компенсує подвійний перелік кожної пари (доданок з~індексами $i,j$
і~доданок з~індексами $j,i$ описують одну й~ту саму пару).

\begin{SCfigure}[1.5][h!]\centering
	\localinput{charges-triangle.tikz}
	\caption{Система трьох точкових зарядів: взаємна енергія набирається із внесків усіх
	пар\label{pic:ChargesTriangle}}
\end{SCfigure}

Формулу~\eqref{eq:PairEnergy} зручно переписати через потенціал. Потенціал поля в~точці
розташування $i$-го заряду, створюваний \emph{усіма іншими} зарядами (крім самого $i$-го,
поле якого в~точці свого ж розташування не визначене),
\begin{equation*}
	\phi_i = \sum\limits_{\substack{j=1 \\ j \neq i}}^n \frac{q_j}{r_{ij}},
\end{equation*}
тому
\begin{equation}\label{eq:SystemEnergyViaPotential}
	\tcbhighmath{
		W = \frac12 \sum\limits_i q_i \phi_i.
	}
\end{equation}

\begin{Attention}
	Формула~\eqref{eq:PairEnergy} описує \emph{лише взаємну} енергію~--- роботу зближення
	\emph{вже готових} зарядів. Власної енергії кожного заряду вона не містить (див.
	\enquote{Власна та взаємна енергія зарядів} нижче).
\end{Attention}


%%% --------------------------------------------------------
\subsection*{Електростатична енергія тіла з~неперервним розподілом заряду}
%%% --------------------------------------------------------




Для тіл з~неперервним розподілом заряду~--- об'ємною густиною $\rho(\vect r)$ та (якщо є)
поверхневою густиною $\sigma(\vect r)$ --- формулу~\eqref{eq:SystemEnergyViaPotential}
узагальнюють природним чином, замінюючи суму на інтеграл, а~$q_i$~--- на елемент заряду
$\rho\, dV$ або $\sigma\, dS$:
\begin{equation}\label{eq:BodyEnergy}
	\tcbhighmath{
		W = \frac12 \iiint\limits_V \rho(\vect r)\phi(\vect r)\, dV + \frac12 \iint\limits_S \sigma(\vect r)\phi(\vect r)\, dS.
	}
\end{equation}

\emph{Приклад 1.} Нехай куля радіуса $R$ несе повний заряд $Q$, рівномірно розподілений
по \emph{поверхні}. Потенціал усюди на поверхні однаковий і дорівнює $\phi = Q/R$, тому
\begin{equation}\label{eq:SphereEnergySurface}
	W = \frac12 \iint\limits_S \sigma\phi\, dS = \frac12\, \sigma\, \frac{Q}{R}\cdot 4\pi R^2 = \frac{Q^2}{2R}.
\end{equation}

\emph{Приклад 2.} Нехай тепер той самий заряд $Q$ розподілений рівномірно по
\emph{об'єму} кулі. Підставивши в~\eqref{eq:BodyEnergy} потенціал рівномірно зарядженої
кулі~\eqref{eq:BallPotential}, отримаємо інший результат (задача~\ref{prb:SphereEnergy}):
\begin{equation}\label{eq:SphereEnergyVolume}
	W = \frac35\, \frac{Q^2}{R}.
\end{equation}

Той самий заряд $Q$ на тілі того самого радіуса дає \emph{різну} енергію залежно від
розподілу заряду (енергії відрізняються в~$6/5$ раза). Це наштовхує на запитання: чи
енергія~--- властивість самих зарядів, чи вона \emph{розподілена в~просторі}, там, де є
поле? Відповідь дасть густина енергії поля.

Застосуємо формулу~\eqref{eq:SphereEnergySurface} до електрона. Якщо уявити його кулькою радіуса $r_e$, зарядженою по
поверхні, і~вважати, що вся енергія спокою $m_ec^2$~--- це енергія поля, то
\begin{equation*}
	\frac{e^2}{2r_e} = m_ec^2 \quad\Longrightarrow\quad r_e = \frac{e^2}{2m_ec^2}\approx1{,}4\cdot10^{-13}~\text{см}.
\end{equation*}
\emph{Класичним радіусом електрона}\index{Радіус електрона класичний} зазвичай називають
$e^2/m_ec^2\approx2{,}8\cdot10^{-13}$~см: множник порядку одиниці залежить від припущення
про розподіл заряду. Але досліди з~розсіяння не виявляють у~електрона жодного розміру аж до
$\sim10^{-16}$~см, а~заряджена кулька мала б розлетітися від власного відштовхування. Тож
класичний радіус~--- не розмір електрона, а~масштаб, на якому енергія поля частинки
порівнянна з~її енергією спокою і~класична електродинаміка перестає працювати.



%% --------------------------------------------------------
\subsection*{Енергія зарядженого конденсатора}
%% --------------------------------------------------------



Розглянемо процес зарядки конденсатора як послідовне перенесення малих порцій заряду $dq$
з~однієї обкладки на іншу (рис.~\ref{pic:CapacitorDq}). Коли на обкладках уже є заряд $q$
(а~різниця потенціалів дорівнює $\Delta\phi = q/C$), перенесення чергової порції $dq$
вимагає роботи $dW = \Delta\phi\, dq = \dfrac{q\,dq}{C}$. Інтегруючи від нульового заряду
до кінцевого значення $q$, знаходимо енергію зарядженого конденсатора:
\begin{equation}\label{eq:CapacitorEnergy}
	\tcbhighmath{
		W = \frac{q^2}{2C} = \frac{C\Delta\phi^2}{2}.
	}
\end{equation}

\begin{SCfigure}[4][h!]\centering
	\localinput{capacitor-dq.tikz}
	\caption{Перенесення заряду $dq$ під час зарядки конденсатора\label{pic:CapacitorDq}}
\end{SCfigure}

\begin{Example}{Числовий приклад} Конденсатор ємністю $C = 1\ \text{мкФ}$, заряджений до напруги $U = 100\ \text{В}$, зберігає заряд
	$q = CU = 10^{-4}\ \text{Кл}$ і~енергію $W = CU^2/2 = 5\ \text{мДж}$~--- зовсім небагато:
	стільки ж енергії потрібно, щоб підняти тіло масою $0{,}5$~кг на $1$~мм. Цінність
	конденсатора в~іншому: цю енергію він здатен віддати \emph{дуже швидко}. Якщо
	розрядити його за $1$~мкс, середня потужність сягне $\sim5$~кВт. Тому конденсатори
	використовують там, де потрібен короткий потужний імпульс: спалахи фотоапаратів,
	імпульсні джерела живлення, дефібрилятори (у~дефібриляторі $C \sim 100$~мкФ заряджають
	до $U \sim 2$~кВ, і~$W = CU^2/2 \sim 200$~Дж віддаються пацієнтові за кілька мілісекунд).
\end{Example}



%% --------------------------------------------------------
\subsection*{Густина енергії електричного поля}\label{sec:FieldEnergyDensity}
%% --------------------------------------------------------



Застосуємо результат~\eqref{eq:CapacitorEnergy} до плоского конденсатора. Оскільки поле
в~ньому однорідне, $\Delta\phi = Ed$, а~ємність дається формулою~\eqref{eq:FlatCap}, маємо
\begin{equation*}
	W = \frac12 C \Delta\phi^2 = \frac12\,\frac{\epsilon S}{4\pi d}\,(Ed)^2 = \frac{\epsilon E^2}{8\pi}\, V,
\end{equation*}
де $V = Sd$~--- об'єм простору між обкладками, тобто об'єм, що його займає поле. Звідси
природно ввести \emphx{об'ємну густину енергії поля}\index{Густина енергії!електричного поля}
\begin{equation}\label{eq:CapacitorEnergyDensity}
	\tcbhighmath{
		w = \frac{\epsilon E^2}{8\pi} = \frac{ED}{8\pi}
	}
\end{equation}
(останню рівність отримано підстановкою $D = \epsilon E$;
див.~розділ~\ref{FieldInDielectrics}).

Ця формула виражає енергію конденсатора не через заряди на обкладках, а~через поле там,
де воно існує. Звідси важлива фізична інтерпретація: носієм енергії є саме \emphx{поле},
і~енергія \emph{локалізована} там, де є поле (у~плоскому конденсаторі~--- лише між
обкладками), а~не \enquote{приписана} зарядам.



%% --------------------------------------------------------
\subsection*{Густина енергії електричного поля (строге виведення)}\label{sec:FieldEnergyDensityStrict}
%% --------------------------------------------------------



Формулу~\eqref{eq:CapacitorEnergyDensity} отримано для окремого випадку однорідного поля,
локалізованого в~проміжку між обкладками плоского конденсатора. Виникає природне
запитання: чи має вона більш загальний характер, чи є наслідком саме цієї спрощеної
геометрії? Покажемо, що вираз для густини енергії електричного поля залишається
справедливим для \emph{довільного} електростатичного поля~--- однорідного чи ні,
локалізованого в~скінченній області чи створеного як завгодно розподіленими зарядами,~---
і~виведемо його строго, виходячи з~енергії взаємодії зарядів.

Нехай заряд системи змінюється на нескінченно малу величину $\delta\rho(\vect r)$.
Оскільки енергія $W=\tfrac12\iiint\rho\phi\,dV$ квадратична за $\rho$ (адже сам потенціал
$\phi$ лінійно виражається через $\rho$), пряме варіювання дає внесок не лише від зміни
$\rho$ при фіксованому $\phi$, а~й~від відповідної зміни самого потенціалу. Записуючи
$\phi\to\phi+\delta\phi$, де $\delta\phi(\vect r)=\iiint \delta\rho(\vect r')/|\vect
r-\vect r'|\,dV'$~--- потенціал, створений приростом заряду,~--- та користуючись симетрією
потенціалу точкового заряду,
\begin{equation*}
	\iiint \rho\,\delta\phi\,dV
	= \iiint\!\!\iiint \frac{\rho(\vect r)\,\delta\rho(\vect r')}{|\vect r-\vect r'|}\,dV\,dV'
	= \iiint \delta\rho\,\phi\,dV,
\end{equation*}
бо відстань між точками не залежить від порядку, $|\vect r-\vect r'|=|\vect r'-\vect r|$.
Отже,
\begin{equation}\label{eq:DeltaW}
	\tcbhighmath{
		\delta W = \iiint\limits_V \delta\rho(\vect r)\, \phi(\vect r)\, dV.
	}
\end{equation}

За теоремою Гаусса для вектора $\Dfield$ у~диференціальній формі, $\delta\rho =
\dfrac1{4\pi}\divg\delta\Dfield$, тому
\begin{equation*}
	\delta W = \frac1{4\pi} \iiint\limits_V \left(\divg \delta\Dfield\right) \phi\, dV.
\end{equation*}
Скористаємось векторною тотожністю $\divg(\phi\, \delta\Dfield) = \phi\,
\divg\delta\Dfield + \delta\Dfield\cdot\vect\nabla\phi$, звідки
\begin{equation*}
	\delta W = \frac1{4\pi} \iiint\limits_V \divg\left(\phi\, \delta\Dfield\right) dV - \frac1{4\pi} \iiint\limits_V \delta\Dfield\cdot\vect\nabla\phi\, dV.
\end{equation*}

Поширимо інтегрування на весь (нескінченний) простір, обравши $V$ як велику кулю радіуса
$R$ із центром десь усередині системи зарядів, і~розглянемо границю $R\to\infty$. Перший
доданок за теоремою Остроградського--Гаусса перетворюється на потік крізь сферу $S_R$:
\begin{equation*}
	\oiint\limits_{S_R} \phi\,\delta\Dfield\cdot d\vect S = \oiint\limits_{S_R} \phi\,\delta D_r\, dS.
\end{equation*}
Для обмеженої системи зарядів на великих відстанях $\phi\sim q/R$, а~$\delta D_r\sim
1/R^2$, тому підінтегральний вираз спадає як $1/R^3$, тоді як площа сфери зростає як
$R^2$. Отже, весь поверхневий інтеграл спадає як $1/R$ і~прямує до нуля при $R\to\infty$:
\begin{equation*}
	\oiint\limits_{S_R} \phi\,\delta D_r\, dS \;\sim\; \frac1R \;\xrightarrow[R\to\infty]{}\; 0.
\end{equation*}
Тому перший доданок зникає, і~залишається лише
\begin{equation*}
	\delta W = -\frac1{4\pi} \iiint\limits_V \delta\Dfield\cdot\vect\nabla\phi\, dV
	= \iiint\limits_V \frac{\Efield\cdot\delta\Dfield}{4\pi}\, dV,
\end{equation*}
де ми скористались зв'язком $\Efield = -\vect\nabla\phi$.

Для лінійного діелектрика ($\Dfield = \epsilon\Efield$) це співвідношення інтегрується
елементарно ($\Efield\cdot\delta\Dfield = \tfrac12\delta(\Efield\cdot\Dfield)$), і~ми
приходимо до загального результату:
\begin{equation}\label{eq:GeneralEnergyDensity}
	\tcbhighmath{
		W = \iiint\limits_V \frac{\Efield\cdot\Dfield}{8\pi}\, dV, \qquad w = \frac{\Efield\cdot\Dfield}{8\pi} = \frac{\epsilon E^2}{8\pi}.
	}
\end{equation}

\begin{Extra}
\emph{Нелінійні середовища*.} Формула~\eqref{eq:GeneralEnergyDensity} спирається на лінійний
зв'язок $\Dfield=\epsilon\Efield$: ним ми скористалися на кроці
$\Efield\cdot\delta\Dfield=\tfrac12\delta(\Efield\cdot\Dfield)$. Проміжний результат
$\delta w = \Efield\cdot\delta\Dfield/4\pi$ загальніший і~справедливий для будь-якого
середовища: формулу~\eqref{eq:DeltaW} тоді слід розуміти як роботу перенесення заряду
$\delta\rho$ з~нескінченності в~точки з~потенціалом $\phi$, а~далі використано лише
скінченність повного заряду та $\Efield=-\vect\nabla\phi$. Тому в~загальному випадку густину
енергії обчислюють як роботу джерел при наростанні поля від нуля:
\begin{equation*}
	w = \frac1{4\pi}\int\limits_0^{\Dfield} \Efield(\Dfield')\cdot d\Dfield',
\end{equation*}
для ізотропного середовища~--- площа під кривою $E(D)$, поділена на $4\pi$. Для середовища
без гістерезису цей інтеграл не залежить від шляху, і~$w(\Dfield)$~--- функція стану. При
$\epsilon=\const$ він дає $w = D^2/8\pi\epsilon = \Efield\cdot\Dfield/8\pi$.
\end{Extra}

\begin{Example}[unbreakable]{Енергія поля точкового заряду}
     Оцінімо енергію поля точкового заряду $q$ у~вакуумі. Густина енергії $w = E^2/(8\pi)
     = q^2/(8\pi r^4)$. Інтегруючи по сферичних оболонках $dV = 4\pi r^2\, dr$ від деякого
     малого радіуса $r_0$ до нескінченності:
	\begin{equation*}
		W = \int\limits_{r_0}^{\infty} \frac{q^2}{8\pi r^4}\cdot 4\pi r^2\, dr = \frac{q^2}{2}\int\limits_{r_0}^{\infty} \frac{dr}{r^2} = \frac{q^2}{2r_0}.
	\end{equation*}
	При $r_0 \to 0$ цей інтеграл \emph{розбігається}~--- власна енергія точкового заряду
	нескінченна. Це~--- не фізична катастрофа, а~сигнал, що \enquote{точковий заряд}~---
	ідеалізація, яка ламається на дуже малих відстанях (там потрібна квантова теорія
	поля). Для реальних заряджених тіл скінченних розмірів (як куля в~прикладах вище)
	власна енергія залишається скінченною.
\end{Example}



%% --------------------------------------------------------
\subsection*{Власна та взаємна енергія зарядів}
%% --------------------------------------------------------



Порівняємо два, на перший погляд, різні вирази для енергії системи зарядів:
\enquote{дискретну} формулу~\eqref{eq:PairEnergy}, $W = \tfrac12\sum q_iq_j/r_{ij}$,
і~\enquote{польову} формулу~\eqref{eq:GeneralEnergyDensity}, $W = \int E^2/(8\pi)\, dV$.
Перша з~них, узагалі кажучи, може набувати як додатних, так і~від'ємних значень
(притягання різнойменних зарядів дає від'ємний внесок); друга ж, як інтеграл від
невід'ємної величини $E^2$,~--- \emph{завжди} невід'ємна. Це здається суперечністю, проте
насправді обидві формули описують не зовсім те саме.

Розглянемо поле системи двох зарядів, $\Efield = \Efield_1 + \Efield_2$. Тоді
\begin{equation*}
	W = \int\limits_V \frac{E^2}{8\pi}\, dV = \int\limits_V \frac{E_1^2}{8\pi}\, dV + \int\limits_V \frac{E_2^2}{8\pi}\, dV + \int\limits_V
	\frac{\Efield_1\cdot\Efield_2}{4\pi}\, dV.
\end{equation*}
Перші два доданки~--- це \emphx{власна енергія}\index{Енергія зарядів!власна} кожного із зарядів окремо; вони завжди
додатні,
\begin{equation*}
	\int\limits_V \frac{E_1^2}{8\pi}\, dV > 0, \qquad \int\limits_V \frac{E_2^2}{8\pi}\, dV > 0,
\end{equation*}
і~не залежать від взаємного розташування зарядів. Третій доданок~--- це \emphx{взаємна
(інтерференційна) енергія}
\begin{equation}\label{eq:SelfMutualEnergy}
	W_\text{вз} = \int\limits_V \frac{\Efield_1\cdot\Efield_2}{4\pi}\, dV,
\end{equation}
яка може мати будь-який знак залежно від взаємного розташування й~знаків зарядів, і~саме
вона точно відповідає \enquote{дискретній} формулі $q_1q_2/r_{12}$ з~попереднього пункту.

Тепер зрозуміло, чому формула~\eqref{eq:PairEnergy} не містить власної енергії точкових
зарядів: поле точкового заряду біля нього необмежено зростає, і~$\int E^2/8\pi\,dV$
розбігається. \enquote{Дискретна} енергія може бути від'ємною, бо вона лише частина
(взаємна) повної, завжди додатної енергії поля, яка містить ще й~нескінченні власні енергії.
Для тіл скінченних розмірів власна енергія скінченна.



%% --------------------------------------------------------
\section{Пондеромоторні сили}\hypertarget{ponderomotive-forces}{}
%% --------------------------------------------------------

Термін \emphx{пондеромоторні сили}\index{Пондеромоторні сили} походить від латинських слів \textit{pondus} (вага,
тягар) та \textit{motor} (той, що рухає). Він виник в~електродинаміці XIX~ст.: так
називали сили, що діють з~боку поля на \enquote{вагомі} (ponderable) тіла~--- провідники,
діелектрики, магнетики,~--- на відміну від \emph{електрорушійних} сил, що діють на самі
заряди (див. історичну довідку в~розділі~\ref{SteadyCurrent}). Сьогодні \emphx{пондеромоторними} називають механічні
сили, з~якими електромагнітне поле діє на макроскопічні тіла.



%% --------------------------------------------------------
\subsection*{Метод віртуальних переміщень}
%% --------------------------------------------------------



Одним з~найефективніших способів обчислення сил, що діють на заряджені провідники
й~діелектрики, є \emphx{метод віртуальних переміщень}\index{Віртуальних переміщень метод}. Його ідея полягає в~тому, щоб
обчислити роботу сил поля $\delta A = F_x\, dx$ при уявному (\emph{віртуальному})
нескінченно малому зміщенні тіла на $dx$ у напрямку, що нас цікавить, після чого шукана
сила знаходиться просто як
\begin{equation*}
	F_x = \frac{\delta A}{dx}.
\end{equation*}

Метод універсальний: він не вимагає з'ясовувати, як саме виникає сила, і~автоматично
враховує всі взаємодії в~системі, бо спирається лише на закон збереження енергії.

Метод спирається на глибший факт~--- \emph{теорему Томсона}: за заданих повних зарядів
провідників рівноважний розподіл заряду мінімізує електростатичну енергію серед усіх
можливих розподілів (за заданих потенціалів мінімізується $W-\sum_i q_i\phi_i$). Якби
інший розподіл давав меншу енергію, заряд самочинно перерозподілився б. Це варіаційний
принцип електростатики, тісно пов'язаний з~теоремою єдиності.



%% --------------------------------------------------------
\subsection*{Сила при сталому заряді провідників}
%% --------------------------------------------------------



Нехай на провідниках підтримується сталий заряд, $q = \const$, тобто в~системі немає
зовнішніх джерел енергії (батарей тощо). За таких умов уся робота $\delta A$ внутрішніх
сил системи при віртуальному переміщенні відбувається виключно за рахунок
\emph{зменшення} власної електричної енергії системи $W$:
\begin{equation*}
	\delta A = -\left(dW\right)_{q=\const}.
\end{equation*}
Якщо шукана сила $F_x$ здійснює роботу $\delta A = F_x\, dx$ при зміщенні тіла на $dx$, то
\begin{equation}\label{eq:ForceConstQ}
	\tcbhighmath{
		F_x = -\left.\frac{dW}{dx}\right|_{q=\const}.
	}
\end{equation}

\begin{Attention}
	Формула~\eqref{eq:ForceConstQ}~--- математичний прийом: у~реальному процесі заряди не
	мусять лишатися сталими. Досить обчислити похідну енергії, \emph{умовно} вважаючи заряд
	фіксованим на тому значенні, яке він має в~даний момент.
\end{Attention}



%% --------------------------------------------------------
\subsection*{Сила при сталій різниці потенціалів}
%% --------------------------------------------------------



Нехай тепер, навпаки, на провідниках підтримуються сталими потенціали, $\Delta\phi =
\const$ (наприклад, провідники приєднані до батареї, яка підживлює заряди в~міру
переміщення тіл і~виконує при цьому роботу). За законом збереження енергії, робота джерела
$\delta A_\text{дж}$ витрачається одночасно на виконання механічної роботи $F_x\, dx$ і~на
зміну енергії поля $dW$:
\begin{equation*}
	\delta A_\text{дж} = F_x\, dx + dW.
\end{equation*}
З~іншого боку, робота джерела з~перенесення заряду $\delta q$ через різницю потенціалів
$\Delta\phi$ дорівнює $\delta A_\text{дж} = \delta q\, \Delta\phi$; оскільки $q =
C\Delta\phi$, а~$\Delta\phi = \const$, маємо $\delta q = \Delta\phi\, \delta C$, тобто
$\delta A_\text{дж} = \Delta\phi^2\, \delta C$. Водночас, за
формулою~\eqref{eq:CapacitorEnergy}, $dW = \tfrac12\Delta\phi^2\, \delta C$ (енергія при
сталому $\Delta\phi$ зростає разом із ємністю). Підставляючи, знаходимо
\begin{equation*}
	F_x\, dx = \Delta\phi^2\,\delta C - \frac12\Delta\phi^2\,\delta C = \frac12\Delta\phi^2\,\delta C = dW,
\end{equation*}
звідки
\begin{equation}\label{eq:ForceConstPhiE}
	\tcbhighmath{
		F_x = +\left.\frac{dW}{dx}\right|_{\Delta\phi=\const}.
	}
\end{equation}

\begin{Attention}[unbreakable]
	Формули~\eqref{eq:ForceConstQ} і~\eqref{eq:ForceConstPhiE} відрізняються \emph{знаком},
	але дають \emph{одну й~ту саму} силу: сила, що діє на тіло в~даний момент, не залежить
	від того, яку уявну зміну системи ми обрали для обчислення. Знак перевертає робота
	джерела, яке при $\Delta\phi=\const$ вкладає в~систему енергію.
\end{Attention}

\begin{SCfigure}[1][tbp]\centering
	\localinput{virtual.tikz}
	\caption{Віртуальне зміщення обкладки: відстань $x$ між обкладками зростає на $dx$, штрихом показано нове положення\label{pic:VirtualDisplacement}}
\end{SCfigure}

\begin{Example}{Притягування обкладок зарядженого конденсатора}

Знайдемо силу взаємного притягання обкладок плоского конденсатора, уявивши, що одна
з~обкладок зміщується так, що відстань $x$ між обкладками зростає на $dx$
(рис.~\ref{pic:VirtualDisplacement}).

\emph{Спосіб 1: $q = \const$.} Зміна енергії поля при зміні відстані між обкладками:
\begin{equation*}
	dW = d\left(\frac{q^2}{2C}\right) = -\frac{q^2}{2C^2}\, dC, \qquad C = \frac{\epsilon S}{4\pi x}, \qquad dC = -\frac{C}{x}\, dx,
\end{equation*}
звідки $dW = \dfrac{q^2}{2Cx}\, dx$. Виражаючи через напруженість поля ($q = CEd = C E x$,
де $x$ тут відіграє роль відстані між обкладками),
\begin{equation*}
	dW = \frac{E^2 x^2\, \epsilon S}{2 \cdot 4\pi x \cdot x}\, dx = \frac{\epsilon E^2 S}{8\pi}\, dx,
\end{equation*}
і~за формулою~\eqref{eq:ForceConstQ}
\begin{equation}\label{eq:ForceCapacitorPlates}
		F_x = -\frac{dW}{dx} = -\frac{\epsilon E^2}{8\pi}\, S.
\end{equation}

\emph{Спосіб 2: $\Delta\phi = \const$.} Аналогічно,
\begin{equation*}
	dW = d\left(\frac{C\Delta\phi^2}{2}\right) = \frac{\Delta\phi^2}{2}\, dC =
	-\frac{\Delta\phi^2 C}{2x}\, dx = -\frac{\epsilon E^2 S}{8\pi}\, dx
\end{equation*}
(тут підставлено $\Delta\phi = Ex$), і~за формулою~\eqref{eq:ForceConstPhiE} знову
отримуємо той самий результат~\eqref{eq:ForceCapacitorPlates}.

Знак \enquote{мінус} у~формулі~\eqref{eq:ForceCapacitorPlates} свідчить про те, що сила
спрямована на \emph{зменшення} відстані між обкладками, тобто обкладки
\emph{притягуються} одна до одної~--- незалежно від знаків зарядів на них, оскільки сила
пропорційна $E^2$.
\end{Example}

Результат~\eqref{eq:ForceCapacitorPlates} загальніший, ніж здається. Тиск поля на будь-яку
заряджену поверхню провідника у~вакуумі (пор. виведення в~розділі~\ref{FieldInMetalls})
\begin{equation*}
	p = \frac{E^2}{8\pi} = 2\pi\sigma^2,
\end{equation*}
де $E = 4\pi\sigma$~--- поле біля поверхні. Цей \emphx{електростатичний
тиск}\index{Електростатичний тиск} завжди напрямлений \emph{назовні} провідника й~діє
в~кожній точці поверхні незалежно від її форми. Помножений на площу обкладки, він дає
силу~\eqref{eq:ForceCapacitorPlates}.

\begin{figure}[tbp]\centering
	\localinput{pull-dielectric.tikz}
	\caption{Конденсатор з~частково всунутою пластиною діелектрика (модель~--- паралельне
		з'єднання вакуумного й~заповненого конденсаторів). На поверхні пластини виступають
		зв'язані заряди; біля краю конденсатора поле неоднорідне, і~його складова $E_\tau$,
		дотична до поверхні пластини, тягне ці заряди всередину~--- звідси сила
		$\vect F$\label{pic:DielectricInCapacitor}}
\end{figure}

\begin{Example}{Сила втягування діелектрика в~конденсатор}
Розглянемо плоский конденсатор, у~який частково (на глибину $x$) всунуто пластину
діелектрика з~проникністю $\epsilon$, і~знайдемо силу, що діє на цю пластину. Зручно
замінити реальну систему еквівалентною~--- двома паралельно з'єднаними конденсаторами:
один, площею $S_1$,~--- вакуумний, а інший, площею $S_2$,~--- повністю заповнений
діелектриком.

За формулою для паралельного з'єднання~\eqref{eq:ParallelCap},
\begin{equation*}
	C = C_1 + C_2 = \frac{S_1}{4\pi d} + \frac{\epsilon S_2}{4\pi d}.
\end{equation*}
Нехай повна довжина конденсатора (у~напрямку висування) дорівнює $L$, ширина обкладок~---
$h$, тоді повна площа $S = Lh$ стала, а $S_1 = \left(1 - \dfrac{x}{L}\right) S$, $S_2 =
\dfrac{x}{L}S$ (тут $x$~--- глибина занурення пластини діелектрика). Підставляючи,
знаходимо ємність як функцію $x$:
\begin{equation*}
	C(x) = \frac{S}{4\pi d}\left(1 + (\epsilon - 1)\frac{x}{L}\right).
\end{equation*}
Обчислимо силу при сталому заряді $q = \const$ за формулою~\eqref{eq:ForceConstQ}:
\begin{equation*}
	F_x = -\frac{dW}{dx} = -\frac{dW}{dC}\frac{dC}{dx}, \qquad \frac{dW}{dC} = -\frac{q^2}{2C^2}, \qquad \frac{dC}{dx} = \frac{S}{4\pi d}\,
	\frac{\epsilon-1}{L},
\end{equation*}
звідки
\begin{equation*}
	F_x = \frac{q^2}{2C^2}\cdot\frac{S}{4\pi d}\cdot\frac{\epsilon - 1}{L} \stackrel{q = CEd}{=} \frac{(\epsilon-1)E^2}{8\pi}\, hd,
\end{equation*}
де ми скористались тим, що в~кожен момент $q = CEd$, а~$S/L = h$. Остаточно,
\begin{equation}\label{eq:ForceDielectric}
		F_x = \frac{(\epsilon-1)E^2}{8\pi}\, S_\text{торця},
\end{equation}
де $S_\text{торця} = hd$~--- площа торцевого перерізу пластини діелектрика (перерізу, яким
вона \enquote{входить} у~конденсатор). Оскільки $\epsilon > 1$, сила $F_x$ додатна~---
вона спрямована в~бік \emph{втягування} діелектрика всередину конденсатора.
\end{Example}

\begin{Attention}
	Метод віртуальних переміщень дає правильну силу, але приховує механізм: сила діє не
	в~області однорідного поля, де на рис.~\ref{pic:DielectricInCapacitor} намальовано
	$\vect F$, а~на краю конденсатора, де поле \emph{неоднорідне} і~на поляризовані
	молекули діє сила $(\vect p\cdot\vect\nabla)\Efield$~\eqref{eq:EforceOnDipole}. Тому
	сила пропорційна $\epsilon - 1$ і~зникає за відсутності поляризації. Так само рідкий
	діелектрик піднімається між зануреними в~нього обкладками (задача~\ref{prb:LiquidDielectricRise}).
\end{Attention}



%% --------------------------------------------------------
\section*{Підсумки}
\phantomsection \addcontentsline{toc}{section}{Підсумки}
%% --------------------------------------------------------

Усі результати цього розділу спираються на одну спільну ідею: електричне поле~--- не
допоміжна розрахункова абстракція, а~фізичний об'єкт, який \emph{сам є носієм енергії}.
Ємність описує, скільки заряду (а~отже, і~поля) система здатна накопичити при заданій
різниці потенціалів; енергія поля визначається інтегралом густини $w =
\Efield\cdot\Dfield/8\pi$ за областю, де це поле існує; а~пондеромоторні сили, що діють на
провідники й~діелектрики, знаходять з~енергії поля методом віртуальних переміщень. Якщо
система ізольована ($q = \const$), сили зменшують її енергію; якщо ж потенціали підтримує
джерело, енергія поля під дією тих самих сил \emph{зростає}~--- за рахунок роботи джерела.

Практично з~цього розділу варто винести три речі. По-перше, ємність будь-якої системи
провідників обчислюється за єдиним алгоритмом: задати заряд, знайти поле, проінтегрувати
різницю потенціалів, поділити. По-друге, густина енергії поля $w =
\Efield\cdot\Dfield/8\pi$~--- формула загальна, а не лише для плоского конденсатора: вона
справедлива для довільного електростатичного поля. По-третє, сили, що діють на провідники
й~діелектрики, найпростіше знаходити методом віртуальних переміщень~--- через похідну
енергії за координатою, з~урахуванням того, яка величина (заряд чи потенціал)
підтримується сталою.



%% --------------------------------------------------------
\Needspace{9\baselineskip}
\section*{Задачі}
\phantomsection \addcontentsline{toc}{section}{Задачі}
%% --------------------------------------------------------

% задачі цього розділу короткі: не розриваємо їх між сторінками

\begin{problem}\label{prb:CapacitanceMatrix}
	\emph{Матриця ємностей*.} Для двох куль із прикладу в~\S~\enquote{Коефіцієнти ємності*} покладіть $R_1=1$~см,
	$R_2=2$~см, $b=50$~см і~обчисліть $C_{11}, C_{22}, C_{12}$ чисельно.
	\begin{solution}
		$C_{11} \approx 1$~см, $C_{22} \approx 2$~см, $C_{12} \approx -0{,}04$~см (знак від'ємний, оскільки $C_{12} \approx -R_1R_2/b$).
	\end{solution}
\end{problem}

\begin{problem}\label{prb:CapacitorBank}
	\emph{Ємність батареї.} Три конденсатори ємностями $C_1 = 1\ \text{мкФ}$, $C_2 = 2\
	\text{мкФ}$, $C_3 = 3\ \text{мкФ}$ з'єднано спочатку послідовно, потім паралельно.
	Обчисліть сумарну ємність у~кожному випадку.
	\begin{solution}
		послідовно $\approx 0{,}55\ \text{мкФ}$; паралельно $6\ \text{мкФ}$.
	\end{solution}
\end{problem}

\begin{problem}\label{prb:UnequalCharges}
	\emph{Заряд однієї обкладки подвоєно.} Плоский конденсатор заряджений звичайним
	чином: $q_1=q$, $q_2=-q$. Заряд однієї з~обкладок збільшили вдвічі, не чіпаючи другої
	(стало $q_1'=2q$, $q_2=-q$). У~скільки разів зміниться напруженість поля \emph{між}
	обкладками? Чи лишається конденсатор при цьому ідеально екранованим? \emph{Вказівка:}
	скористайтеся розкладом~\eqref{eq:UnequalPlatesField}.
	\begin{solution}
		поле між обкладками зросте в~$1{,}5$ раза; екранування порушиться~--- з'явиться поле й~назовні, $E_\text{зовні}=2\pi q/S$.
	\end{solution}
\end{problem}

\begin{problem}\label{prb:SphereEnergy}
	\emph{Енергія поля зарядженої кулі.} Обчисліть енергію електричного поля рівномірно
	зарядженої по об'єму кулі радіуса $R$ із повним зарядом $Q$ двома способами: (а) через
	формулу~\eqref{eq:BodyEnergy}, знайшовши потенціал усередині кулі; (б) через інтеграл
	$\int E^2/(8\pi)\, dV$ за всім простором. Переконайтесь, що результати збігаються. \emph{Вказівка:} поле
	всередині кулі $E = Qr/R^3$, зовні~--- $E = Q/r^2$.
	\begin{solution}
		$W=\dfrac{3Q^2}{5R}$ обома способами.
	\end{solution}
\end{problem}

\begin{problem}\label{prb:SphericalCapEnergy}
	\emph{Енергія сферичного конденсатора.} Вакуумний сферичний конденсатор з~радіусами
	обкладок $R_1<R_2$ заряджено зарядом $q$. Знайдіть енергію поля інтегруванням густини
	$w=E^2/8\pi$ і~порівняйте з~$q^2/2C$. Яка частка енергії зосереджена в~шарі
	$R_1<r<2R_1$, якщо $R_2\gg R_1$? Який тиск поле чинить на внутрішню обкладку?
	\begin{solution}
		$W=\displaystyle\int_{R_1}^{R_2}\frac{q^2}{8\pi r^4}\,4\pi r^2dr=\frac{q^2}{2}\left(\frac1{R_1}-\frac1{R_2}\right)=\frac{q^2}{2C}$; при $R_2\gg R_1$ у~шарі $R_1<r<2R_1$ зосереджена половина енергії; тиск $p=E^2/8\pi=q^2/8\pi R_1^4$, напрямлений назовні.
	\end{solution}
\end{problem}

\begin{problem}\label{prb:SlabBattery}
	\emph{Діелектрик із батареєю і~без неї.} У~плоский вакуумний конденсатор ємністю $C$,
	заряджений до напруги $U$, повністю вставляють пластину діелектрика з~проникністю
	$\epsilon$, що заповнює весь проміжок. Як зміняться заряд, напруга й~енергія
	конденсатора, якщо (а)~конденсатор перед цим від'єднали від батареї; (б)~батарея
	лишається приєднаною? Яку роботу виконує поле над пластиною і~яку~--- батарея?
	Втягується пластина чи виштовхується?
	\begin{solution}
		(а)~$q$ стале, $U'=U/\epsilon$, $W'=W/\epsilon$, де $W=CU^2/2$; поле виконує над пластиною роботу $W(1-1/\epsilon)$ за рахунок своєї енергії. (б)~$U$ стала, $q'=\epsilon q$, $W'=\epsilon W$; батарея переносить заряд $(\epsilon-1)CU$ і~виконує роботу $(\epsilon-1)CU^2=2\Delta W$, половина якої йде на приріст енергії поля, а~половина $(\epsilon-1)CU^2/2$~--- на роботу над пластиною. В~обох випадках пластина втягується, хоча енергія поля в~(а) зменшується, а~в~(б) зростає.
	\end{solution}
\end{problem}

\begin{problem}\label{prb:ForceConstPhi}
	\emph{Сила при сталій напрузі.} Плоский конденсатор із площею обкладок $S=20\
	\text{см}^2$ приєднано до батареї, що підтримує напругу $U=300$~В; відстань між
	обкладками $d=0{,}5$~мм. Обчисліть силу притягання обкладок безпосередньо за
	формулою~\eqref{eq:ForceConstPhiE} (через залежність $C(x)$), не користуючись готовим
	результатом~\eqref{eq:ForceCapacitorPlates}, і~переконайтесь, що відповідь збігається.
	\emph{Вказівка:} у~СІ $C=\epsilon_0S/x$.
	\begin{solution}
		$F=\epsilon_0SU^2/2d^2\approx 3{,}2\times10^{-3}\ \text{Н}$.
	\end{solution}
\end{problem}

\begin{problem}\label{prb:DielectricForce}
	\emph{Сила, що діє на діелектрик.} У~плоский конденсатор з~обкладками шириною $h = 5\
	\text{см}$ і~відстанню між ними $d = 2\ \text{мм}$ частково всунуто пластину
	діелектрика з~$\epsilon = 4$. Поле в~конденсаторі $E = 5\times 10^4\ \text{В/м}$.
	Обчисліть силу, що втягує пластину всередину конденсатора. \emph{Вказівка:} у~СІ
	формула~\eqref{eq:ForceDielectric} має вигляд $F =
	\dfrac{(\epsilon-1)\epsilon_0E^2}{2}\,hd$.
	\begin{solution}
		$F \approx 3{,}3\times 10^{-6}\ \text{Н}$.
	\end{solution}
\end{problem}

\begin{problem}\label{prb:LiquidDielectricRise}
	\emph{Втягування рідкого діелектрика.} Плоский конденсатор із вертикально
	розташованими обкладками (площа кожної \(S = l\cdot h_0\), відстань між обкладками
	\(d\)) з'єднано з~джерелом постійної напруги \(U\). Нижні кінці обкладок занурено
	в~рідину з~діелектричною проникністю \(\epsilon\) і~густиною \(\rho_\text{р}\), яка не
	проводить струм. Рідина втягується в~проміжок між обкладками і~встановлюється на
	деякій висоті \(h\) (відносно рівня рідини поза конденсатором; \(h<h_0\), крайовими
	ефектами й~капілярністю знехтуйте).

	Знайдіть:
	\begin{enumerate}
		\item висоту підйому рідини \(h\);
		\item кількість теплоти \(Q\), що виділилася при встановленні рівноваги.
	\end{enumerate}

	\emph{Вказівка:} скористайтесь тим, що при сталій напрузі втягувальна сила
	\(F_{\text{ел}} = \dfrac{(\epsilon-1)l}{4\pi d}\dfrac{U^2}{2}\) (це
	формула~\eqref{eq:ForceDielectric} при \(E=U/d\), \(S_\text{торця}=ld\)) не залежить
	від глибини занурення, а~зміна енергії конденсатора \(\Delta W\) і~робота джерела
	\(A_{\text{дж}}\) пов'язані співвідношенням \(A_{\text{дж}} = 2\Delta W\). При
	рівновазі втягувальна сила зрівноважується вагою піднятого стовпа рідини. Теплоту
	знайдіть із закону збереження енергії \(A_{\text{дж}}=\Delta W+\tfrac12 mgh+Q\), де
	\(\tfrac12 mgh\)~--- приріст потенціальної енергії піднятої рідини.
	\begin{solution}
		\(h = \dfrac{(\epsilon-1)U^2}{8\pi\rho_\text{р} g d^2}\); \quad \(Q = \dfrac{1}{2}mgh = \dfrac{(\epsilon-1)^2 l U^4}{128\pi^2\rho_\text{р} g d^3}\),
		де \(m=\rho_\text{р} l d h\)~--- маса піднятого стовпа.
	\end{solution}
\end{problem}

\begin{Questions}
\begin{quizlist}
	\item Що називають ємністю відокремленого провідника? Чому відношення $C = q/\phi$ не
	  залежить від заряду і~від чого воно залежить насправді?
	\item Чому означення $C = q/\phi$ втрачає сенс, якщо поблизу є інші провідники? Як
	  тоді визначається взаємна ємність пари провідників?
	\item Що таке конденсатор і~чому його взаємна ємність майже не залежить від навколишніх
	  тіл?
	\item Чому дорівнює ємність двох далеких одна від одної куль
	  (рис.~\ref{pic:TwoSpheres}) у~головному наближенні? Що описує поправка, пропорційна
	  $1/b$?
	\item Виведіть формули для ємності батареї при послідовному та паралельному з'єднанні
	  конденсаторів (рис.~\ref{pic:CapacitorsConnection}). Що залишається однаковим
	  у~кожному з~випадків: заряд чи напруга?
	\item Чому послідовне з'єднання завжди зменшує сумарну ємність, а~паралельне збільшує?
	  Яка величина для конденсатора відіграє роль опору?
	\item Сформулюйте триетапний алгоритм обчислення ємності. Застосуйте його для плоского
	  конденсатора та запишіть результат.
	\item Виведіть ємність сферичного конденсатора. Що станеться з~формулою при
	  $R_2\to\infty$ і~чому цей результат очікуваний?
	\item Яка загальна структура формул для ємності плоского, сферичного та циліндричного
	  конденсаторів? Що в~них відіграє роль \enquote{геометричного розміру}?
	\item Чому дорівнює взаємна енергія системи точкових зарядів? Навіщо в~формулі $W =
	  \tfrac12\sum\sum q_iq_j/r_{ij}$ множник $\tfrac12$ і~чому вона не враховує власну
	  енергію зарядів?
	\item Запишіть енергію системи зарядів через потенціали. Чому енергія кулі, зарядженої
	  по поверхні, відрізняється від енергії кулі того ж заряду й~радіуса, зарядженої по
	  об'єму? Яке запитання це ставить про локалізацію енергії?
	\item Виведіть енергію зарядженого конденсатора $W = q^2/2C = C\Delta\phi^2/2$ через
	  роботу перенесення заряду $dq$.
	\item Як з~енергії плоского конденсатора отримують густину енергії поля $w = \epsilon
	  E^2/8\pi = ED/8\pi$? Чому цю формулу можна вважати загальною, а~не лише для
	  однорідного поля?
	\item Яку роль відіграє те, що потік вектора $\phi\,\delta\Dfield$ крізь нескінченно
	  віддалену поверхню дорівнює нулю, у~виведенні загального виразу для густини енергії?
	\item Чому енергія \enquote{дискретної} формули для точкових зарядів може бути
	  від'ємною, а~енергія, обчислена як інтеграл від $E^2/8\pi$, завжди додатна? Що таке
	  власна та взаємна (інтерференційна) енергія?
	\item Чому власна енергія точкового заряду нескінченна? Що це означає для моделі
	  \enquote{точковий заряд}?
	\item Що таке пондеромоторні сили та звідки походить термін? У~чому полягає метод
	  віртуальних переміщень?
	\item Виведіть формули для сили при $q=\const$ та при $\Delta\phi=\const$. Чому вони
	  відрізняються знаком, хоча описують одну й~ту саму фізичну силу?
	\item Діелектричну пластину вставляють у~конденсатор, від'єднаний від батареї і~приєднаний до неї. Чому в~першому випадку енергія поля зменшується, а~в~другому зростає, хоча пластина втягується в~обох?
	\item Знайдіть силу притягання обкладок плоского конденсатора. Чому вона напрямлена на
	  зменшення відстані між обкладками незалежно від знаків зарядів?
	\item Чому діелектрик втягується в~область сильного поля конденсатора? Як у~формулі
	  для сили з'являється множник $(\epsilon-1)$ і~що з~цього випливає для вакууму?
\end{quizlist}
\end{Questions}

\Closesolutionfile{answer}
